\begin{figure}[t]
\centering
% Shared color TikZ styles for the paper figures.
\colorlet{paperBlue}{blue!9!white}
\colorlet{paperBlueLine}{blue!55!black}
\colorlet{paperTeal}{teal!10!white}
\colorlet{paperTealLine}{teal!55!black}
\colorlet{paperAmber}{orange!15!white}
\colorlet{paperAmberLine}{orange!60!black}
\colorlet{paperRose}{red!9!white}
\colorlet{paperRoseLine}{red!55!black}
\colorlet{paperViolet}{violet!10!white}
\colorlet{paperVioletLine}{violet!55!black}
\colorlet{paperInk}{black!72}

\tikzset{
  paperfig/.style={
    font=\small,
    >=Stealth,
    line cap=round,
    line join=round
  },
  paperline/.style={
    -{Stealth[length=2.1mm,width=1.6mm]},
    line width=0.58pt,
    draw=paperInk
  },
  paperdash/.style={
    paperline,
    dashed,
    draw=black!48
  },
  paperbox/.style={
    draw=black!58,
    line width=0.58pt,
    rounded corners=2pt,
    fill=white,
    align=center,
    inner xsep=5pt,
    inner ysep=5pt,
    minimum height=8mm
  },
  paperdata/.style={
    paperbox,
    fill=paperBlue,
    draw=paperBlueLine
  },
  paperproc/.style={
    paperbox,
    fill=paperTeal,
    draw=paperTealLine
  },
  papermodel/.style={
    paperbox,
    fill=paperAmber,
    draw=paperAmberLine,
    line width=0.65pt
  },
  paperout/.style={
    paperbox,
    fill=paperRose,
    draw=paperRoseLine,
    line width=0.65pt,
    font=\small\bfseries
  },
  paperstate/.style={
    paperbox,
    fill=paperViolet,
    draw=paperVioletLine
  },
  papernote/.style={
    font=\scriptsize,
    align=center,
    text=black!70,
    fill=white,
    inner sep=1.4pt
  },
  paperaxis/.style={
    line width=0.45pt,
    draw=black!72
  },
  papergrid/.style={
    line width=0.25pt,
    draw=black!15
  },
  paperplot/.style={
    line width=0.95pt,
    draw=paperBlueLine
  },
  paperplotA/.style={
    line width=0.95pt,
    draw=paperBlueLine
  },
  paperplotB/.style={
    line width=0.95pt,
    draw=paperAmberLine
  },
  papermark/.style={
    circle,
    draw=paperBlueLine,
    fill=white,
    line width=0.55pt,
    inner sep=1.6pt
  },
  papermarkA/.style={
    papermark,
    draw=paperBlueLine,
    fill=paperBlue
  },
  papermarkB/.style={
    papermark,
    draw=paperAmberLine,
    fill=paperAmber
  }
}

\resizebox{0.98\linewidth}{!}{%
\begin{tikzpicture}[paperfig, x=1mm, y=1mm]
  \node[paperdata, text width=18mm] (task) at (7,34) {任务输入\\\scriptsize 题面/签名/测试};
  \node[papermodel, text width=20mm] (gen) at (31,34) {单步候选\\\scriptsize 一次生成};
  \node[paperproc, text width=17mm] (exec) at (51,34) {执行\\判定};

  \draw[paperline] (task) -- (gen);
  \draw[paperline] (gen) -- (exec);

  \node[paperout, text width=16mm] (done) at (51,18) {通过\\记录结果};
  \draw[paperline] (exec) -- node[papernote, right] {pass} (done);

  \node[paperstate, text width=19mm] (residual) at (31,18) {残差样本\\\scriptsize 单步失败};
  \draw[paperline] (exec) -- node[papernote, above] {fail} (residual);

  \node[paperproc, text width=17mm] (repair) at (7,1) {修补\\\scriptsize 错误回注};
  \node[paperproc, text width=17mm] (route) at (23,1) {分流\\\scriptsize 难度门控};
  \node[paperproc, text width=17mm] (rescue) at (39,1) {救援\\\scriptsize 重链闭环};
  \node[paperproc, text width=17mm] (cascade) at (55,1) {级联\\\scriptsize 候选覆盖};

  \draw[paperline] (residual) -- (repair);
  \draw[paperline] (repair) -- (route);
  \draw[paperline] (route) -- (rescue);
  \draw[paperline] (rescue) -- (cascade);
  \draw[paperline] (cascade.east) to[bend right=28] node[papernote, right] {筛选} (done.east);

  \draw[paperAmberLine, line width=0.45pt] (4,-9) -- (58,-9);
  \foreach \x/\lab/\txt in {
    7/单步/最低成本,
    17/修补/反馈预算,
    31/分流/结构预算,
    45/救援/失败预算,
    55/级联/覆盖预算
  } {
    \draw[paperAmberLine, line width=0.45pt] (\x,-10.5) -- (\x,-7.5);
    \node[font=\scriptsize, align=center] at (\x,-14) {\lab\\\txt};
  }

  \draw[paperTealLine, line width=0.7pt] (4,7) -- ++(0,3) -- ++(8,0) -- ++(0,3);
  \draw[paperTealLine, line width=0.7pt] (20,7) -- ++(0,3) -- ++(8,0) -- ++(0,3);
  \draw[paperTealLine, line width=0.7pt] (36,7) -- ++(0,3) -- ++(8,0) -- ++(0,3);
  \node[papernote, text width=46mm] at (31,10)
    {新增预算只作用于失败后的残差样本，逐级增加反馈、路由和候选覆盖};
\end{tikzpicture}%
}
\caption{预算分档编排整体框架}
\label{fig:overall_framework}
\end{figure}
